%------------------------------------------------------------------------------%
\begin{question}
  Escreva a equação
  \[
    x^2 + xy + y^2 = 1
  \]
  em coordenadas polares e simplifique
\end{question}

%------------------------------------------------------------------------------%
\begin{resolution}{Solução}
  \onslide<+->{Sabemos que
  \[
    x = r\cos(\theta)
    \qquad\qquad
    y = r\sin(\theta)
    \qquad\qquad
    r^2 = x^2 + y^2
  \]}
  \vspace{-\baselineskip}
  \begin{align*}
    \onslide<+->{x^2 + xy + y^2                                         & = 1 \\}
    \onslide<+->{x^2 + y^2 + xy                                         & = 1 \\}
    \onslide<+->{r^2 + r\cos\theta r\sin\theta                          & = 1 \\}
    \onslide<+->{r^2 + r^2\cos\theta \sin\theta                         & = 1 \\}
    \onslide<+->{r^2\left(1 + \cos\theta \sin\theta \right)             & = 1 &}
    \onslide<+->{r & = \frac{\pm 1}{\sqrt{1 + \cos\theta \sin\theta}}}
  \end{align*}
\end{resolution}
%------------------------------------------------------------------------------%
\begin{resolution}{Solução Alternativa}
  \begin{align*}
    \onslide<+->{r^2\left(1 + \cos\theta \sin\theta \right)   & = 1 \\}
    \onslide<+->{r^2\left(1 + \frac{\sin(2\theta)}{2} \right) & = 1 \\}
    \onslide<+->{r^2\left(2 + \sin(2\theta) \right)           & = 2 \\}
    \onslide<+->{r & = \frac{\pm 2}{\sqrt{2 + \sin(2\theta)}}}
  \end{align*}
\end{resolution}

%------------------------------------------------------------------------------%
